\section{Round-eleven positive closure: sliced-current filtering with zero-evidence compactification}
\label{sec:r11-c1}

The observation state is not the space of arbitrary measures on level sets.
It is the graph closure of finite-mass normal currents for which Federer slices
are defined, together with a projective direction at zero evidence.

\subsection{The stratified slicing complex}
Let $\mathcal N_W^r(S)$ denote weighted normal $r$-currents on a smooth
stratum, with mass and boundary mass.  A regular observation map
$O^u:S\to\mathbb R^d$ acts by the slice
$\langle T,O^u,y\rangle$ for almost every $y$, not by ordinary restriction
to the level set.

\begin{theorem}[Closed observation--collision current tower]
\label{thm:r11-c1-slice}
For a finite family of stratified submersive observations, prediction by the
B2 trace evolution and observation by Federer slicing generate a countable
locally finite current complex.  The slice operators satisfy the coarea mass
identity, preserve orientations and collision involutions, and are closed in
the weighted mass-plus-boundary norm.  Critical sets are represented by their
own lower-rank strata with induced slice currents.
\end{theorem}

\begin{proof}
Start with smooth densities and collision traces.  The slicing theorem for
normal currents gives, for every Lipschitz observation chart,
\[
 \int \mathbf M(\langle T,O,y\rangle)dy
 \le \operatorname{Lip}(O)^d\mathbf M(T).
\]
Apply the boundary-slice identity for normal mass.  Refine the Whitney
stratification after every finite chronological block; local finiteness follows
from the B2 singularity/trace complexity bound.  Closure of normal-current
slices defines the update on the completed domain, and prediction preserves
that domain by the Green graph theorem.
\end{proof}

\subsection{Evidence and its boundary}
For a nonzero unnormalized current $\Lambda$, write
$Z=\langle1,\Lambda\rangle$, $z=\log Z$, and
$[\Lambda]=\Lambda/Z$.  Compactify by adjoining pairs
$(0,-\infty,\vartheta)$, where $\vartheta$ is a projective limit direction
in the compact weighted current sphere.  The direction records which
normalized posterior can arise as evidence vanishes.

\begin{theorem}[Feller evidence kernel]
\label{thm:r11-c1-feller}
The prediction--slice--normalization update extends to a Feller kernel on the
zero-evidence compactification.  At positive evidence it is the ordinary Bayes
posterior.  At the boundary, transition probabilities are weighted by the
unnormalized evidence and depend continuously on the retained projective
direction; observations of exactly zero evidence carry zero probability.
\end{theorem}

\begin{proof}
Theorem~\ref{thm:r11-c1-slice} gives continuity of the unnormalized sliced
current.  Polar decomposition in the weighted current cone is continuous
away from zero and compact after projectivization.  Any sequence with
$Z_n\downarrow0$ has a projective subsequence; adjoining that direction makes
the state compact.  In the integrated Bayes kernel, the factor $Z_n$
multiplies every future payoff, so different zero-mass directions affect only
the explicitly retained boundary coordinate and do not create a discontinuity
in the unnormalized transition.
\end{proof}

\subsection{Observation coefficients from the trajectory expansion}
The admissible observation is a finite vector of additive particle-path and
actual-contact cylinders.  It may be global in space, but it is additive over
the exact marked trajectory configuration and therefore enters the B2 source
before clustering.

\begin{theorem}[Source/control-uniform inserted coefficient]
\label{thm:r11-c1-coefficient}
On every regular observation stratum and compact control/source chart, the
joint constraint--observation characteristic function has Gaussian behavior
near its exact saddle, exponential contraction on every compact nonzero
annulus, and speed-dependent decay at high frequency.  With a bounded central
insertion $F$, numerator and denominator have the same local coefficient
operator, and their ratio converges uniformly to the conditional inserted
expectation.
\end{theorem}

\begin{proof}
Adjoin the observation vector to the B2 marked source and to the B1 dynamic
block mark.  The affine submersion condition on each regular stratum gives the
extra rank; B2 polymer decoupling gives linearly many conditional smoothing
blocks.  The proof of the B1 mixed Cramer theorem applies to the enlarged
vector.  A central insertion changes the leading left/right eigenvectors but
not the saddle exponent; differentiating the inserted block operator gives
its explicit coefficient.  Fourier inversion and division yield the ratio.
\end{proof}

\subsection{Reachable correlations and control}
Define the reachable analytic belief class by a fixed complex neighborhood of
the complete factorial-current generating functional, a uniform radius, and a
finite declared list of order-one phase coordinates.  This is an a priori
normed class, not ``whatever correlations survive.''

\begin{theorem}[Quenched reachable-chaos invariance]
\label{thm:r11-c1-chaos}
Prediction, every regular positive-evidence observation update, and the
zero-evidence boundary kernel preserve the reachable analytic belief class.
Outside a compact subchart, the evidence carries the explicit large-deviation
cost.  Higher unlisted cumulants remain $O(C^j\mu^{1-j})$ uniformly over all
admissible strategies.
\end{theorem}

\begin{proof}
Normal convergence of the B2 trajectory source on a common complex ball gives
Cauchy estimates for all factorial-current cumulants.  The inserted coefficient
of Theorem~\ref{thm:r11-c1-coefficient} divides two analytic functions only on
positive-evidence compact charts, preserving the radius.  Boundary directions
are retained rather than normalized.  Induction over decision blocks and the
strategy-uniform source bounds give the cumulant estimate.
\end{proof}

\begin{theorem}[Belief game, reduction, and Bernstein--von Mises]
\label{thm:r11-c1-game}
The full belief state has a well-defined dynamic program.  On the reachable
analytic class, projection to the declared density/contact/phase coordinates
changes each transition kernel and reward by $o(1)$ in a weighted Wasserstein
metric; a Bellman Lipschitz recursion therefore gives a uniform $o(1)$ game
value error.  For an identifiable finite-dimensional regular phase parameter,
a prior with positive continuous density, and the inserted LAN expansion from
Theorem~\ref{thm:r11-c1-coefficient}, the posterior satisfies a canonical
Bernstein--von Mises theorem away from phase boundaries.
\end{theorem}

\begin{proof}
Theorem~\ref{thm:r11-c1-feller} and compact control sets give the full dynamic
program.  The cumulant bound in Theorem~\ref{thm:r11-c1-chaos} gives a coupling
of full and reduced transitions with one-step error $o(1)$; backward Bellman
iteration sums those errors over the fixed horizon.  For the statistical
statement, expand the exact finite log likelihood at local alternatives using
the inserted coefficient.  Identifiability gives a positive information
matrix, prior regularity controls the denominator, and exponential posterior
tightness removes phase-boundary neighborhoods.
\end{proof}

\begin{theorem}[Round-eleven C1 closure]
\label{thm:r11-c1-main}
The controlled exact-observation model has a rigorously sliced normal-current
state, a zero-evidence projective compactification, an inserted local
coefficient theorem for additive path/contact observations, a noncircular
reachable analytic class, and a stable belief-game and posterior limit.
\end{theorem}

\begin{proof}
Combine Theorems~\ref{thm:r11-c1-slice},
\ref{thm:r11-c1-feller}, \ref{thm:r11-c1-coefficient},
\ref{thm:r11-c1-chaos}, and \ref{thm:r11-c1-game}.
\end{proof}
